\documentclass[10pt,a4paper]{amsart}
\usepackage{amsmath,amssymb,amsthm,graphicx}
% Retained minimal preamble for the Stacks Project chapter sites-cohomology.tex (Cohomology on Sites).
% Verbatim excerpt of the chapter preamble: Xy-pic setup (native preamble lines 14-18) and the native
% theorem declarations with the mathematical macros (native preamble lines 154-217). Copied mechanically;
% nothing added, removed or redefined. The wrapper class and the amsmath/amssymb/amsthm/graphicx base
% packages are supplied by the assembler, so the class/externaldocument part of the native preamble is omitted.
% For commutative diagrams we use Xy-pic
\usepackage[all]{xy}

% We use 2cell for 2-commutative diagrams.
\xyoption{2cell}
% Theorem environments.
%
\theoremstyle{plain}
\newtheorem{theorem}[subsection]{Theorem}
\newtheorem{proposition}[subsection]{Proposition}
\newtheorem{lemma}[subsection]{Lemma}

\theoremstyle{definition}
\newtheorem{definition}[subsection]{Definition}
\newtheorem{example}[subsection]{Example}
\newtheorem{exercise}[subsection]{Exercise}
\newtheorem{situation}[subsection]{Situation}

\theoremstyle{remark}
\newtheorem{remark}[subsection]{Remark}
\newtheorem{remarks}[subsection]{Remarks}

\numberwithin{equation}{subsection}

% Macros
%
\def\lim{\mathop{\mathrm{lim}}\nolimits}
\def\colim{\mathop{\mathrm{colim}}\nolimits}
\def\Spec{\mathop{\mathrm{Spec}}}
\def\Hom{\mathop{\mathrm{Hom}}\nolimits}
\def\Ext{\mathop{\mathrm{Ext}}\nolimits}
\def\SheafHom{\mathop{\mathcal{H}\!\mathit{om}}\nolimits}
\def\SheafExt{\mathop{\mathcal{E}\!\mathit{xt}}\nolimits}
\def\Sch{\mathit{Sch}}
\def\Mor{\mathop{\mathrm{Mor}}\nolimits}
\def\Ob{\mathop{\mathrm{Ob}}\nolimits}
\def\Sh{\mathop{\mathit{Sh}}\nolimits}
\def\NL{\mathop{N\!L}\nolimits}
\def\CH{\mathop{\mathrm{CH}}\nolimits}
\def\proetale{{pro\text{-}\acute{e}tale}}
\def\etale{{\acute{e}tale}}
\def\QCoh{\mathit{QCoh}}
\def\Ker{\mathop{\mathrm{Ker}}}
\def\Im{\mathop{\mathrm{Im}}}
\def\Coker{\mathop{\mathrm{Coker}}}
\def\Coim{\mathop{\mathrm{Coim}}}

% Boxtimes
%
\DeclareMathSymbol{\boxtimes}{\mathbin}{AMSa}{"02}

%
% Macros for moduli stacks/spaces
%
\def\QCohstack{\mathcal{QC}\!\mathit{oh}}
\def\Cohstack{\mathcal{C}\!\mathit{oh}}
\def\Spacesstack{\mathcal{S}\!\mathit{paces}}
\def\Quotfunctor{\mathrm{Quot}}
\def\Hilbfunctor{\mathrm{Hilb}}
\def\Curvesstack{\mathcal{C}\!\mathit{urves}}
\def\Polarizedstack{\mathcal{P}\!\mathit{olarized}}
\def\Complexesstack{\mathcal{C}\!\mathit{omplexes}}
% \Pic is the operator that assigns to X its picard group, usage \Pic(X)
% \Picardstack_{X/B} denotes the Picard stack of X over B
% \Picardfunctor_{X/B} denotes the Picard functor of X over B
\def\Pic{\mathop{\mathrm{Pic}}\nolimits}
\def\Picardstack{\mathcal{P}\!\mathit{ic}}
\def\Picardfunctor{\mathrm{Pic}}
\def\Deformationcategory{\mathcal{D}\!\mathit{ef}}
\usepackage{cleveref}
\crefname{originallemma}{original lemma}{original lemmas}
\Crefname{originallemma}{Original lemma}{Original lemmas}
\makeatletter
\newlabel{lemma-derived-limit-is-ok}{{0941}{source}{}{sourcepointer-1}{}}
\newlabel{lemma-derived-limit-is-ok@cref}{{[originallemma][0][]0941}{[1][0][]source}{}{}{}}
\makeatother
\crefname{originallemma}{original lemma}{original lemmas}
\Crefname{originallemma}{Original lemma}{Original lemmas}
\makeatletter
\newlabel{injectives-lemma-derived-products}{{07D9}{source}{}{sourcepointer-2}{}}
\newlabel{injectives-lemma-derived-products@cref}{{[originallemma][0][]07D9}{[1][0][]source}{}{}{}}
\makeatother
\crefname{originallemma}{original lemma}{original lemmas}
\Crefname{originallemma}{Original lemma}{Original lemmas}
\makeatletter
\newlabel{lemma-adjoint}{{07A6}{source}{}{sourcepointer-3}{}}
\newlabel{lemma-adjoint@cref}{{[originallemma][0][]07A6}{[1][0][]source}{}{}{}}
\makeatother
\crefname{originallemma}{original lemma}{original lemmas}
\Crefname{originallemma}{Original lemma}{Original lemmas}
\makeatletter
\newlabel{categories-lemma-adjoint-exact}{{0038}{source}{}{sourcepointer-4}{}}
\newlabel{categories-lemma-adjoint-exact@cref}{{[originallemma][0][]0038}{[1][0][]source}{}{}{}}
\makeatother
\title{The derived pushforward commutes with derived limits}
\author{The Stacks Project Authors}
\date{stacks-project:sites-cohomology.tex:local-snapshot-20261011; GNU FDL 1.2 or later}
\begin{document}
\maketitle
\noindent\textit{Editorial note.} All delivered mathematical text is the author's own, copied from the pinned Stacks Project chapter file sites-cohomology.tex; no statement or proof step was rewritten, reordered or removed.\par\medskip
\noindent\textit{Editorial note.} The proof is the author's complete argument; the author additionally supplies an alternative proof in the case of enough points, and both are inside the counted span.\par\medskip
\noindent\textit{Editorial note.} Every labelled reference occurring inside the counted spans is carried in assembly.external\_references as a non-load-bearing navigation pointer (same-chapter ones at their definition site, cross-chapter ones at the invocation inside the counted span), so no load-bearing mathematics is replaced by a stub. No bibliography citation occurs inside the counted spans.\par\medskip
\noindent\textit{Editorial note.} The checked chapter list for same-name collisions is the set of stacks chapters staged in this repository: chow, cohomology, divisors, etale-cohomology, flat, more-algebra, more-morphisms, morphisms, spaces-more-morphisms, varieties together with this card's chapter sites-cohomology.\par\medskip
\noindent\textit{Editorial note.} The preamble is the chapter's own retained minimal preamble (sources/stacks-sites-cohomology/retained-preamble.tex), a verbatim excerpt of the native preamble: the Xy-pic setup and the native theorem declarations with the mathematical macros.\par\medskip
\noindent\textit{Allowed standard prerequisites (editorial).}\begin{itemize}
\item Derived pushforward along a morphism of ringed topoi and its computation by K-injective representatives. Same chapter, the construction of Rf\_* and its computation using K-injective complexes. Supplies the functor whose commutation with derived limits is asserted and the representatives used in the proof.
\item Derived limit R lim of an inverse system of complexes and its defining triangle. Same chapter together with the chapter Derived Categories, the construction of R lim through the product and the shift, and the associated distinguished triangle. Supplies the object R lim K\_n and the triangle used in both proofs.
\item Commutation of the underived pushforward with products and of restriction with limits. Same chapter, the behaviour of f\_* with respect to products of sheaves and the exactness properties used in the alternate proof. Supplies the comparison making both proofs work.
\end{itemize}\medskip
\ifdefined\hypertarget\hypertarget{sourcepointer-1}{}\fi
\noindent\textit{Source reference.} Original lemma 0941: sources/stacks-sites-cohomology/pinned-original-source.tex, source lines 5067--5067. Navigation pointer invoked at proof line 5174; the counted span contains the pointer only and the cited result is not reproduced inside it, so the pointer is registered as a non-load-bearing navigation stub in the established form for Stacks material. No bibliography citation is involved.\par\medskip
\ifdefined\hypertarget\hypertarget{sourcepointer-2}{}\fi
\noindent\textit{Source reference.} Original lemma 07D9: sources/stacks-sites-cohomology/pinned-original-source.tex, source lines 5119--5119. Navigation pointer invoked at proof line 5119; the counted span contains the pointer only and the cited result is not reproduced inside it, so the pointer is registered as a non-load-bearing navigation stub in the established form for Stacks material. No bibliography citation is involved.\par\medskip
\ifdefined\hypertarget\hypertarget{sourcepointer-3}{}\fi
\noindent\textit{Source reference.} Original lemma 07A6: sources/stacks-sites-cohomology/pinned-original-source.tex, source lines 3990--3990. Navigation pointer invoked at proof line 5192; the counted span contains the pointer only and the cited result is not reproduced inside it, so the pointer is registered as a non-load-bearing navigation stub in the established form for Stacks material. No bibliography citation is involved.\par\medskip
\ifdefined\hypertarget\hypertarget{sourcepointer-4}{}\fi
\noindent\textit{Source reference.} Original lemma 0038: sources/stacks-sites-cohomology/pinned-original-source.tex, source lines 5194--5194. Navigation pointer invoked at proof line 5194; the counted span contains the pointer only and the cited result is not reproduced inside it, so the pointer is registered as a non-load-bearing navigation stub in the established form for Stacks material. No bibliography citation is involved.\par\medskip
\begin{lemma}
\label{lemma-Rf-commutes-with-Rlim}
Let $f : (\Sh(\mathcal{C}), \mathcal{O}) \to (\Sh(\mathcal{C}'), \mathcal{O}')$
be a morphism of ringed topoi. Then $Rf_*$ commutes with $R\lim$, i.e.,
$Rf_*$ commutes with derived limits.
\end{lemma}

\begin{proof}
Let $(K_n)$ be an inverse system of objects of $D(\mathcal{O})$.
By induction on $n$ we may choose actual complexes $\mathcal{K}_n^\bullet$
of $\mathcal{O}$-modules and maps of complexes
$\mathcal{K}_{n + 1}^\bullet \to \mathcal{K}_n^\bullet$ representing the
maps $K_{n + 1} \to K_n$ in $D(\mathcal{O})$. In other words, there exists
an object $K$ in $D(\mathcal{C} \times \mathbf{N})$ whose associated inverse
system is the given one. Next, consider the commutative diagram
$$
\xymatrix{
\Sh(\mathcal{C} \times \mathbf{N}) \ar[r]_g \ar[d]_{f \times 1} &
\Sh(\mathcal{C}) \ar[d]_f \\
\Sh(\mathcal{C}' \times \mathbf{N}) \ar[r]^{g'} &
\Sh(\mathcal{C}')
}
$$
of morphisms of topoi. It follows that
$R\lim R(f \times 1)_*K = Rf_* R\lim K$. Working through the definitions
and using Lemma \ref{lemma-derived-limit-is-ok}
we obtain that $R\lim (Rf_*K_n) = Rf_*(R\lim K_n)$.

\medskip\noindent
Alternate proof in case $\mathcal{C}$ has enough points. Consider the defining
distinguished triangle
$$
R\lim K_n \to \prod K_n \to \prod K_n
$$
in $D(\mathcal{O})$. Applying the exact functor $Rf_*$ we obtain
the distinguished triangle
$$
Rf_*(R\lim K_n) \to Rf_*\left(\prod K_n\right) \to Rf_*\left(\prod K_n\right)
$$
in $D(\mathcal{O}')$. Thus we see that it suffices to prove that
$Rf_*$ commutes with products in the derived category (which are not just
given by products of complexes, see
Injectives, Lemma \ref{injectives-lemma-derived-products}).
However, since $Rf_*$ is a right adjoint by Lemma \ref{lemma-adjoint}
this follows formally (see
Categories, Lemma \ref{categories-lemma-adjoint-exact}).
Caution: Note that we cannot apply
Categories, Lemma \ref{categories-lemma-adjoint-exact}
directly as $R\lim K_n$ is not a limit in $D(\mathcal{O})$.
\end{proof}

\end{document}
